
\subsection*{Abstract}
Production Knowledge Graph reasoning systems require more than structural traversal --- they must handle time-varying facts, propagate uncertainty through multi-hop paths, accommodate nodes that belong to multiple communities simultaneously, and support continuous improvement of their core attention parameters. We present CEREBRUM's \textbf{Algorithmic Depth} layer (Phase 17), five orthogonal enhancements to the core CSA reasoning engine that collectively enable temporal, probabilistic, and adaptive reasoning without introducing training data requirements or sacrificing the zero-hallucination guarantee. The five components are: (1) temporal edge validity windows with decay; (2) uncertainty propagation through the CSA formula; (3) soft community membership with fractional overlap scores; (4) \texttt{CSAParameterLearner} --- online, training-free CSA weight adaptation from query feedback; and (5) KGE integration (TransE \cite{bordes2013transe} / RotatE \cite{sun2019rotate}) as optional drop-in embedding providers. Each component is independently composable; the full suite achieves +14.2\% relative H@10 on MetaQA-3hop over the Phase 16 baseline. In v2.71.0, the CSA formula has been expanded to 10 parameters and the learning stack upgraded through Phases 45--48: parameter persistence, auto-retrain scheduling, and adaptive search strategy further extend the algorithmic depth concept to runtime adaptation.

\subsection*{1. Introduction}
The core CSA formula [Buchorn, 2026] was designed with algebraic simplicity as a primary constraint: six weighted terms, a sigmoid activation, and configurable per-community parameter overrides. This design deliberately excludes temporal dynamics, uncertainty semantics, and continuous learning to ensure mathematical transparency. However, real-world KG deployments exhibit all three: facts have validity periods, sources have varying reliability, and query traffic provides a continuous signal about which reasoning strategies are working.

Phase 17 adds five capabilities as composable layers that augment the core without modifying it, preserving backward compatibility and the mathematical interpretability of every reasoning step.

\subsection*{2. Temporal Edge Validity}

\subsubsection*{2.1 Edge Temporal Metadata}
Each edge is extended with optional temporal metadata:
\begin{lstlisting}
Edge.valid_from: Optional[float]    # Unix timestamp (start of validity)
Edge.valid_until: Optional[float]   # Unix timestamp (end of validity)
Edge.temporal_weight: float = 1.0  # Current weight after decay
\end{lstlisting}

\subsubsection*{2.2 Temporal Decay Function}
For edges with a \texttt{valid\_until} timestamp, temporal weight decays exponentially after the validity period ends:

$$w_{temp}(t) = w_0 \cdot \exp\left(-\lambda \cdot \max(0, t - t_{until})\right)$$

The decay constant $\lambda$ is configurable per relation type (e.g., \texttt{CURRENT\_PRICE} decays faster than \texttt{BORN\_IN}). Edges with \texttt{temporal\_weight < \textbackslash\{\}epsilon\_\{min\}} (default: 0.01) are automatically removed by the REM Cycle.

\subsubsection*{2.3 Integration with CSA}
The temporal weight multiplicatively modulates the CSA attention weight:

$$a_{temp}(u,v,k) = a(u,v,k) \cdot w_{temp}(t_{query})$$

where $t_{query}$ is the snapshot time at query start (consistent with Query Snapshot Isolation, [Buchorn, 2026]).

\subsection*{3. Uncertainty Propagation}

\subsubsection*{3.1 Per-Edge Confidence Scores}
Edges ingested via \texttt{IngestionPipeline} carry a \texttt{confidence} attribute (default: 1.0). This confidence represents source reliability at ingest time and is stored as edge metadata.

\subsubsection*{3.2 Path-Level Uncertainty}
For a reasoning path $P = \{e_1, e_2, \ldots, e_L\}$, the path confidence propagates as:

$$\text{conf}(P) = \prod_{i=1}^{L} c_i^{\alpha} \cdot \left(1 - \beta \cdot \text{Var}(\{c_i\})\right)$$

where $c_i$ is the confidence of edge $e_i$, $\alpha$ controls sensitivity to low-confidence edges, and the variance term penalizes paths with inconsistent confidence profiles (a path mixing one very-high and one very-low confidence edge is less trustworthy than a path with uniformly moderate confidence).

\subsubsection*{3.3 Uncertainty in Answer Extraction}
The \texttt{AnswerExtractor} appends per-path uncertainty bounds to the output:

\begin{lstlisting}[language=Json]
{
    "answer": "Marie Curie",
    "path": ["Physics", "Nobel_Prize_1903", "Marie_Curie"],
    "csa_score": 0.847,
    "path_confidence": 0.763,
    "confidence_interval": [0.71, 0.81]
}
\end{lstlisting}

\subsection*{4. Soft Community Membership}

\subsubsection*{4.1 Motivation}
Hard community assignment (each node belongs to exactly one community) is appropriate for highly modular graphs but produces sharp discontinuities at community boundaries. Nodes on community boundaries --- particularly Hub nodes with connections to multiple clusters --- receive CSA penalties that systematically under-weight their structural importance.

\subsubsection*{4.2 Fractional Membership Scores}
Soft membership extends the community map to store a probability distribution over communities for each node:

$$\mu_v = \{c_1: p_1, c_2: p_2, \ldots, c_K: p_K\}, \quad \sum_k p_k = 1$$

The primary community is $\arg\max_k p_k$. The secondary membership is used by the community consensus term $S_C(u,v)$:

$$S_C^{soft}(u,v) = \sum_{k} p_k^{(u)} \cdot p_k^{(v)}$$

This is the dot product of the two membership distributions, which equals 1 for perfectly same-community nodes, 0 for fully disjoint membership, and a continuous value in between for partially-overlapping nodes.

\subsubsection*{4.3 Computation}
Soft membership scores are derived from the DSCF modularity matrix: the raw modularity contributions $\Delta Q_{vk}$ for each node-community pair are normalized via softmax with a temperature parameter $\tau$ (default: 2.0). Higher $\tau$ produces softer (more uniform) distributions; lower $\tau$ approaches hard assignment.

\subsection*{5. CSAParameterLearner}

\subsubsection*{5.1 The Learning Problem}
The six CSA weights $(\alpha, \beta, \gamma, \delta, \varepsilon, \zeta)$ are initialized from theoretical priors (semantic similarity and community structure should dominate, with relation weight and decay as secondary terms). However, different graph domains have different optimal weightings: causal graphs may benefit from higher $\gamma$ (relation weight); temporal graphs may need higher $\delta$ (recency penalty); social graphs may need higher $\beta$ (community consensus).

\subsubsection*{5.2 Online Gradient-Free Adaptation}
The \texttt{CSAParameterLearner} adapts weights using a query-feedback signal without gradient computation:

\textbf{Feedback signal}: After a query completes, if the user explicitly validates or rejects the top answer, the learner records a $+1$ or $-1$ signal along with the attention weight breakdown for the winning path.

\textbf{Update rule} (coordinate-wise moving average):
$$\theta_i^{(t+1)} = (1 - \eta) \cdot \theta_i^{(t)} + \eta \cdot r_t \cdot a_i^{(t)}$$

where $r_t \in \{+1, -1\}$ is the feedback signal, $a_i^{(t)}$ is the contribution of weight $i$ to the winning path score, and $\eta$ is the learning rate (default: 0.05).

\textbf{Constraints}: All weights are projected back to the simplex $\sum_i \theta_i = 1, \theta_i \geq 0.01$ after each update, ensuring no term is completely suppressed.

\subsubsection*{5.3 Per-Community Parameters}
The \texttt{CSAParameterLearner} maintains separate parameter sets per community (consistent with Community-Specific CSA Parameters, [Buchorn, 2026]), enabling different communities to learn different optimal weightings from the same query traffic.

\subsection*{6. KGE Embedding Integration}

\subsubsection*{6.1 TransE}
TransE [Bordes et al., 2013] models each relation $r$ as a translation vector: a triple $(h, r, t)$ is valid iff $\vec{h} + \vec{r} \approx \vec{t}$. CEREBRUM integrates TransE as an optional drop-in for the \texttt{EmbeddingEngine}:

\begin{lstlisting}[language=Python]
kge = TransEEngine(adapter, dim=128, margin=1.0, epochs=100)
kge.train()
embedding_engine = KGEEmbeddingAdapter(kge)
\end{lstlisting}

\subsubsection*{6.2 RotatE}
RotatE [Sun et al., 2019] models relations as rotations in complex embedding space, handling symmetric, antisymmetric, inverse, and compositional relations. Its embeddings provide better CSA semantic similarity scores for graphs with rich relational diversity.

\subsubsection*{6.3 Integration with CSA}
KGE embeddings are used exclusively in the $\cos(\vec{e}_u, \vec{e}_v)$ term of the CSA formula. All other terms (community structure, relation weight, distance penalty, hop decay, PageRank) continue to use graph-structural features. This hybrid design preserves the interpretability of the non-embedding terms while upgrading the semantic similarity signal.

\subsection*{7. Recent Advances (v2.51.1 -> v2.71.0)}

\subsubsection*{7.1 10-Parameter CSA Formula (Phase 43/45)}
The original 6-parameter CSA formula has been expanded to a 10-parameter formulation:

$$a(u,v,k) = \text{sigmoid}(\alpha \cdot sim + \beta \cdot cs + \gamma \cdot etw - \delta \cdot nd + \varepsilon \cdot hd + \zeta \cdot pr_v + \eta \cdot td + \iota \cdot nr_v - \mu \cdot sd + \theta \cdot grounding)$$

The four new parameters are:
\begin{itemize}
  \item $\zeta$ (\textbf{PageRank prior}): Boosts structurally central destination nodes
  \item $\eta$ (\textbf{Temporal decay}): Penalizes edges with expired validity windows
  \item $\iota$ (\textbf{Node recency}): Rewards recently-updated or recently-traversed nodes
  \item $\mu$ (\textbf{Synthesis-density penalty}): Penalizes paths over-relying on Synaptic Bridge-synthesized edges (REM Engine)
  \item $\theta$ (\textbf{Grounding confidence}): Rewards edges with high ingest-time confidence scores
\end{itemize}


Default weights: \texttt{(0.4, 0.4, 0.1, 0.05, 0.05, 0.1, 0.1, 0.05, 0.1, 1.0)}

The \texttt{CSAParameterLearner} was correspondingly upgraded from 5-parameter to \textbf{10-parameter} operation (Phase 45), with correct penalty signs ($-\delta$, $-\mu$) applied during SGD gradient steps.

\subsubsection*{7.2 MetaParameterLearner Online SGD (Phase 22/45)}
The \textbf{MetaParameterLearner} operates on all 10 features with correct penalty sign handling. It receives \texttt{POST /feedback} signals and updates community-specific parameter vectors via online SGD with a configurable learning rate. The \texttt{MetaParameterLearner.to\_dict()} / \texttt{from\_dict()} methods enable checkpoint serialization.

\subsubsection*{7.3 Parameter Persistence (Phase 47)}
\texttt{MetaParameterLearner} checkpoints can be exported via \texttt{GET /params}, stored as JSON, and restored via \texttt{POST /params} or the \texttt{--params-file} CLI flag at startup. This enables warm-start deployments where a previously trained parameter set is loaded before the first query, eliminating cold-start variance.

\subsubsection*{7.4 Auto-Retrain Scheduler (Phase 48)}
The \texttt{CSAParameterLearner.fit()} method (batch gradient descent over accumulated \texttt{(pos, neg)} path pairs) is now triggered automatically via \texttt{POST /retrain}. The auto-retrain scheduler fires when the feedback buffer exceeds a configurable threshold (default: 100 pairs), running \texttt{fit()} in a background thread without blocking query processing. After each retrain, the global prior is updated and community-specific parameters are re-initialized from the new prior.

\subsubsection*{7.5 Adaptive Search Strategy (Phase 53)}
Phase 53 extends the "algorithmic depth" concept to \textbf{runtime structural adaptation}: rather than using fixed beam parameters, the traversal engine measures local graph density at query time and selects beam width, depth limit, and branching factor dynamically. This is the first instance of CEREBRUM's reasoning parameters being driven by graph structure at inference time rather than at configuration time, completing the arc from static parameters (Phase 17) to online learned parameters (Phase 22/45/48) to dynamically adapted parameters (Phase 53).

\subsection*{8. Prior Art Differentiation}

\textbf{Temporal edges vs. temporal KG systems:} TNTComplEx \cite{lacroix2020tntcomplex}, TTransE \cite{lin2015ttranse}, and HyTE \cite{sun2017hyte} embed entity-time pairs in a joint space, requiring timestamped training data. CEREBRUM's temporal decay is a parameter applied to edge metadata at query time --- purely structural, no training required.

\textbf{Uncertainty propagation vs. probabilistic KG systems:} ProBase \cite{wu2012probase} and NELL \cite{carlson2010nell} maintain confidence scores but do not propagate uncertainty through multi-hop paths. CEREBRUM's variance-penalized path confidence is computed analytically per-query.

\textbf{Soft community vs. overlapping community detection:} BIGCLAM \cite{yang2013bigclam} and DEMON \cite{coscia2012demon} detect overlapping communities but produce static memberships offline. CEREBRUM's soft membership is derived from the live DSCF modularity matrix, updating automatically after each rebalance.

\textbf{CSAParameterLearner vs. meta-learning:} MAML \cite{finn2017maml} and Reptile \cite{nichol2018reptile} require gradient computation over a differentiable loss. \texttt{CSAParameterLearner} uses coordinate-wise moving averages over a binary feedback signal --- no gradients, no backpropagation, no training data requirement.

\textbf{KGE integration vs. pure embedding methods:} KGQA systems like EmbedKGQA \cite{saxena2020improve} use KGE embeddings as the primary reasoning mechanism. CEREBRUM uses them as one of ten terms in the CSA formula, where graph topology, community structure, and PageRank continue to dominate the reasoning signal.

\subsection*{9. Experimental Results}

Combined Phase 17 enhancement suite evaluated on MetaQA (zero-shot, full-graph):

\begin{table}[h]
\small\centering
\begin{tabular}{llll}
\toprule
Configuration & H@10 (1-hop) & H@10 (3-hop) & $\Delta$ vs. Phase 16 \\
\midrule
Phase 16 baseline & 0.960 & 0.248 & --- \\
+ Temporal edges & 0.971 & 0.329 & +3.5% \\
+ Uncertainty propagation & 0.960 & 0.337 & +6.0% \\
+ Soft community & 0.972 & 0.348 & +9.4% \\
+ CSAParameterLearner & 0.974 & 0.353 & +11.0% \\
+ KGE (RotatE) embeddings & 0.976 & 0.363 & +14.2% \\
\bottomrule
\end{tabular}
\end{table}


All five components compose independently and additively.

\textbf{v2.71.0 canonical benchmark results} (full 10-parameter CSA, MetaParameterLearner, adaptive search):

\begin{table}[h]
\small\centering
\begin{tabular}{lll}
\toprule
Benchmark & Metric & v2.71.0 Result \\
\midrule
MetaQA 1-hop & H@1 / H@10 & 83.2% / 99.0% \\
MetaQA 2-hop & H@1 / H@10 & 63.3% / 94.3% \\
MetaQA 3-hop & H@1 / H@10 & 56.8% / 90.7% \\
WebQSP OPT & H@1 / H@10 / MRR & 6.27% / 20.84% / 10.66% \\
IKGWQ & AUC & 0.89 \\
GrailQA & F1 / H@1 & 19.6% / 13.0% \\
\bottomrule
\end{tabular}
\end{table}


\subsection*{10. Conclusion}
The Algorithmic Depth layer demonstrates that meaningful reasoning improvements can be achieved through principled, composable algorithmic extensions rather than increased model size or training data. The five Phase 17 components collectively advance H@10 by 14.2\% on the hardest benchmark while preserving complete interpretability of every reasoning step. In v2.71.0, the evolution continues: the 10-parameter CSA formula, online MetaParameterLearner, parameter persistence, auto-retrain scheduling, and adaptive search strategy extend the algorithmic depth concept from static composition to a fully adaptive reasoning pipeline that improves automatically with usage.

\textit{See also:} \textbf{Paper 022} --- Looped Beam Traversal (Phase 70) adds a further dimension of algorithmic depth: iterative refinement via LoopLM-style looping [zhu2025loooplm]. Rather than a fixed traversal depth H, the number of reasoning \textit{passes} T becomes an adaptive parameter governed by an exit gate. This is the natural successor to the algorithmic depth concept --- depth is now dynamic in two dimensions: hop depth (H) within a pass, and pass count (T) across iterations.